% Part: propositional-logic
% Chapter: propositional-logic 
% Section: completeness

\documentclass[../../../include/open-logic-section]{subfiles}

\begin{document}

\olfileid{pl}{prp}{com}

\olsection{ప్రతిపాదనా తర్కం యొక్క సంపూర్ణత}

\begin{defn}
\ollabel{def:MaxCon}
\tetoken{సూత్రాల}{!!{formula}s} సమితి $\Gamma$ సంగతమైనదై, $\Gamma \subseteq \Delta$ అయ్యే ఏ సంగత సమితి $\Delta$కైనా $\Gamma = \Delta$ అయితే, దానిని \emph{గరిష్ఠ సంగత సమితి} అంటారు.
\end{defn}

\begin{lem}[సత్య ఉపసిద్ధాంతం]
  \ollabel{lem:truth}
$\Gamma$ గరిష్ఠ సంగత సమితి అనుకుందాం. అప్పుడు:
\begin{enumerate}
\item \ollabel{lem:truth-1} $!A \in \Gamma$ అవడం $\lnot !A \notin \Gamma$ అవడానికి అవసరమైనదీ చాలినదీ;
\item \ollabel{lem:truth-2} $\Gamma \Proves !A$ అవడం $!A \in \Gamma$ అవడానికి అవసరమైనదీ చాలినదీ;
\item \ollabel{lem:truth-3} $!A \lif !B \in \Gamma$ అవడం $!A \notin \Gamma$ లేదా $!B \in \Gamma$ అవడానికి అవసరమైనదీ చాలినదీ.
\end{enumerate}
\end{lem}

\begin{proof}
\begin{enumerate}
\item $!A \in \Gamma$ అనుకుందాం. $\lnot !A \in \Gamma$ కూడా అయితే, $\Gamma$ అసంగతమైనది. $!A$, $\lnot!A$ రెండూ $\Gamma$లో లేకపోతే, \olref[axd]{prop:phi} ప్రకారం $\Gamma\cup\{!A\}$, $\Gamma \cup \{\lnot !A\}$లలో కనీసం ఒకటి సంగతమైనది; అప్పుడు $\Gamma$ గరిష్ఠమైనది కాదు. వ్యతిరేక దిశను కూడా ఇదే విధంగా నిరూపించవచ్చు.
\item $!A \in \Gamma$ అయితే, నిశ్చయంగా $\Gamma \Proves !A$. మరోవైపు $\Gamma \Proves !A$ అనుకుందాం. $!A \notin \Gamma$ అయితే, \olref{lem:truth-1} ప్రకారం $\lnot !A \in \Gamma$. కాబట్టి $\Gamma \Proves !A$, $\Gamma \Proves \lnot !A$ రెండూ వస్తాయి; ఇది వైరుధ్యం.
\item అభ్యాసం.
\end{enumerate}
\end{proof}

\olref{lem:truth}లోని \olref{lem:truth-2} యొక్క ఎడమ నుంచి కుడికి దిశ $\Gamma$ \emph{నిగమనాత్మకంగా సంవృతమైనదని} చూపిస్తుంది; అంటే అన్ని~$!A$లకూ $\Gamma \Proves !A$ అయితే $!A \in \Gamma$.

\begin{prob}
\olref{lem:truth} యొక్క నిరూపణను పూర్తి చేయండి.
\end{prob}

\begin{thm}[సంపూర్ణత]
\ollabel{thm:completeness}
$\Gamma$ సంగతమైనదైతే, $\Gamma$ సంతృప్తిపరచదగినది.
\end{thm}

\begin{proof}
$!A_0$, $!A_1$,~\dots అనేవి భాషలోని అన్ని \tetoken{సూత్రాల}{!!{formula}s} సమగ్ర జాబితా అనుకుందాం. \tetoken{సూత్రాల}{!!{formula}s} సమితుల ఆరోహణ శ్రేణి $\Gamma_0$, $\Gamma_1$,~\dotsను పునరావృత్తంగా ఇలా నిర్వచిస్తాం:
\begin{align*}
\Gamma_0 & = \Gamma\\
\Gamma_{n+1} &  =
\begin{cases}
  \Gamma_n \cup \{!A_n\} & \text{$\Gamma_n \cup \{!A_n\}$ సంగతమైనదైతే;}\\
  \Gamma_n \cup \{\lnot !A_n\} & \text{లేకపోతే}.
\end{cases}
\end{align*}
తర్వాత ఇలా నిర్వచిస్తాం:
\[
\Gamma^* = \bigcup_{0\le n}\Gamma_n.
\]
నిరూపణలో కింది విషయాలను వరుసగా స్థాపిస్తాం:
\begin{enumerate}
\item ప్రతి $n$కు $\Gamma_n$ సంగతమైనది (\olref[axd]{prop:phi}ను ఉపయోగించి $n$పై ఆగమనంతో);
\item $\Gamma^*$ సంగతమైనది;
\item $\Gamma^*$ గరిష్ఠమైనది.
\end{enumerate}
ఇప్పుడు $v(p_i) = \True$ అవడం $p_i \in \Gamma^*$ అవడానికి అవసరమైనదీ చాలినదీ అయ్యేలా విలువ నిర్ధారణ $v$ను నిర్వచిస్తాం. $!A$పై ఆగమనంతో, మరింత సంక్లిష్ట \tetoken{వాక్యాలకు}{!!{sentence}s} కూడా $\Gamma^*$లో సభ్యత్వం $v$ ప్రకారం సత్యత్వంతో సమానమని చూపిస్తాం: $\overline{v}(!A) = \True$ అవడం $!A \in \Gamma^*$ అవడానికి అవసరమైనదీ చాలినదీ. ముఖ్యంగా $v$ అనేది $\Gamma^*$ను సంతృప్తిపరుస్తుంది. $\Gamma \subseteq \Gamma^*$ కాబట్టి, కోరుకున్నట్లు $\Gamma$ కూడా సంతృప్తిపరచదగినది.
\end{proof}

\begin{cor}
$\Gamma \Entails !A$ అయితే $\Gamma \Proves !A$.
\end{cor}

\begin{proof}
$\Gamma\Proves/ !A$ అయితే, \olref[axd]{prop:prov-incons} ప్రకారం $\Gamma \cup \{\lnot !A\}$ సంగతమైనది. సంపూర్ణత సిద్ధాంతం ప్రకారం $\Gamma \cup \{\lnot !A\}$ సంతృప్తిపరచదగినది; అందువల్ల $\Gamma \Entails/ !A$.
\end{proof}

\begin{prop}[సంహతత సిద్ధాంతం]
\ollabel{prop:compactness}
$\Gamma$ సంతృప్తిపరచదగినది అవడం, $\Gamma$ యొక్క ప్రతి \emph{పరిమిత} ఉపసమితి $\Gamma_0$ సంతృప్తిపరచదగినది అవడానికి అవసరమైనదీ చాలినదీ.
\end{prop}

\begin{proof}
$\Gamma$ సంతృప్తిపరచలేనిది అవడం అది అసంగతమైనది అవడానికి అవసరమైనదీ చాలినదీ; ఇది $\Gamma$ యొక్క ఏదో పరిమిత ఉపసమితి $\Gamma_0$ అసంగతమైనది అవడానికి అవసరమైనదీ చాలినదీ; ఇది $\Gamma$ యొక్క ఏదో పరిమిత ఉపసమితి $\Gamma_0$ సంతృప్తిపరచలేనిది అవడానికి అవసరమైనదీ చాలినదీ.
\end{proof}

\begin{cor}
$\Gamma \Entails !A$ అవడం, $\Gamma$ యొక్క ఏదో పరిమిత ఉపసమితి $\Gamma_0$కు $\Gamma_0 \Entails !A$ అవడానికి అవసరమైనదీ చాలినదీ.
\end{cor}

\end{document}
